\documentclass[11pt]{article}
\usepackage[utf8]{inputenc}
\usepackage{amsmath, amssymb}
\usepackage{hyperref}
\usepackage{booktabs}
\usepackage{tabularx}
\usepackage{newunicodechar}

% Ensure the Lean arrow character compiles correctly in pdfLaTeX
\newunicodechar{→}{\ensuremath{\to}}

\title{Formal Foundations of Systems Engineering: A Comparative Analysis of Set-Theoretic and Dependent Type-Theoretic Paradigms}
\author{}
\date{}


\usepackage{csquotes}
\usepackage[style=numeric,backend=biber]{biblatex}
\addbibresource{refs.bib}

\begin{document}
\maketitle

\section{Introduction: The Epistemological Crisis in Systems Engineering}

The discipline of systems engineering is currently navigating a profound epistemological transition. As the exponential proliferation of cyber-physical systems, autonomous networks, and socio-technical infrastructures accelerates, the historical reliance on document-centric, heuristic methodologies has proven fundamentally inadequate. Driven by the necessity to manage immense complexity, the field is actively migrating toward Model-Based Systems Engineering (MBSE) \parencite{Ref1, Ref2}. However, this transition has exposed a critical vulnerability: the prevailing methodologies and graphical modeling languages utilized in MBSE, such as the Systems Modeling Language (SysML), largely lack rigorous mathematical underpinnings \parencite{Ref3, Ref4}. While established engineering disciplines such as aerospace, civil, and mechanical engineering rely heavily on the robust, deterministic mathematical foundations of calculus, differential equations, and finite element analysis, the architectural and behavioral models of systems engineering remain largely informal, descriptive, and disconnected from computable mathematics \parencite{Ref3}. This foundational deficit is not merely an academic concern; it poses severe risks to the development of safety-critical systems \parencite{Ref5, Ref6}. Ambiguous requirements and informal specifications frequently lead to semantic misinterpretations during the transition from system architecture to component implementation, resulting in catastrophic failures, immense cost overruns, and prolonged development schedules \parencite{Ref2, Ref7}. The International Council on Systems Engineering (INCOSE) has formally recognized this crisis through its Future of Systems Engineering (FuSE) initiative \parencite{Ref8, Ref9}. Conceived to realize the SE Vision 2035, the FuSE program established a dedicated ``Foundations'' stream tasked with assessing the adequacy of the theoretical bases of the discipline \parencite{Ref8, Ref10}. The explicit mandate of this stream is to determine future directions by rooting systems engineering in rigorous mathematical theory, information theory, and the sciences of complexity, elevating the discipline from a practice of heuristics to a mathematically formalized science \parencite{Ref8, Ref10}. Historically, the most rigorous and exhaustive attempt to provide a mathematical foundation for systems engineering was A. Wayne Wymore's Tricotyledon Theory of System Design (T3SD) \parencite{Ref11, Ref12}. Developed in the late twentieth century, T3SD is a mathematically comprehensive framework grounded entirely in classical set theory. It formalized the conceptualization of system states, inputs, outputs, transition functions, and coupling recipes, providing a theoretically sound blueprint for system specification and architecture \parencite{Ref11, Ref13, Ref14}. However, because T3SD relies on classical Zermelo-Fraenkel set theory, its implementation in contemporary, automated digital environments has been severely constrained by the inherent difficulties of computing and mechanizing set-theoretic constructs \parencite{Ref15, Ref16}. Set theory, designed as a foundation for classical mathematics, is inherently descriptive rather than prescriptive, making the automated verification of complex system behaviors computationally intractable \parencite{Ref17}. In parallel with the evolution of systems engineering, the fields of computer science and formal logic experienced a paradigm shift with the maturation of Dependent Type Theory (DTT) \parencite{Ref18, Ref19}. Formulations such as Martin-Löf Type Theory and the Calculus of Constructions have emerged as powerful alternative foundations for mathematics \parencite{Ref18, Ref20}. DTT serves as the logical bedrock for modern interactive theorem provers and proof assistants \parencite{Ref19, Ref21, Ref22}. By integrating logic and computation natively through the Curry-Howard isomorphism, DTT offers a paradigm where system specifications are mathematically represented as types, and the physical or software systems that satisfy those specifications are represented as computable programs or terms \parencite{Ref22}. This report presents an exhaustive comparative analysis of the set-theoretic approach (typified by Wymore's T3SD) and the dependent type-theoretic approach to modeling discrete and continuous systems. Utilizing a direct encoding of Wymorian discrete systems into a modern proof assistant as an empirical context for correctness and provability \parencite{Ref23, Ref23}, this report explores the deep philosophical, logical, and practical differences between set theory and type theory. Furthermore, it details how translating systems engineering frameworks into DTT conceptually elevates MBSE, enabling correct-by-construction architectures, eliminating the need for vast libraries of manual relational proofs, and directly supporting the foundational mandates of the INCOSE FuSE initiative. Finally, this report proposes targeted research trajectories tailored for premier academic journals, emphasizing the profound methodological benefits of type-theoretic frameworks to the broader systems engineering community.

\section{Literature Review: The Foundational Schism Between Set Theory and Type Theory}

To fully contextualize the methodological divergence between Wymorian T3SD and modern formalized systems engineering, it is strictly necessary to examine the underlying mathematical foundations. The schism between classical set theory and dependent type theory is not merely a matter of syntactic preference or notation; it represents fundamentally divergent conceptions of logic, equality, function spaces, and computational behavior \parencite{Ref18, Ref23, Ref24, Ref25}. This section reviews the literature defining these differences and articulates what each paradigm enables or prevents within the context of mathematical modeling.

\subsection{Classical Set Theory: Material Ontology and External Logic}
Classical set theory, primarily formalized as Zermelo-Fraenkel set theory with the Axiom of Choice (ZFC), serves as the traditional foundation of mathematics \parencite{Ref16}. It is the language in which nearly all twentieth-century engineering mathematics, including Wymore's T3SD, is expressed \parencite{Ref11}. The ontology of classical set theory is fundamentally ``material'' in nature \parencite{Ref18}. In a material framework, mathematical elements exist independently of and prior to the sets that contain them. Consequently, it is mathematically permissible and logically valid to ask whether an arbitrary, undefined entity $A$ is an element of an arbitrary set $X$ (denoted as $A \in X$) \parencite{Ref18}. Because elements are not intrinsically categorized, set theory relies entirely on an external logic---specifically, first-order predicate logic \parencite{Ref23}. The primitive logical notions of truth, falsity, conjunction, disjunction, implication, and universal or existential quantification exist completely outside the universe of sets \parencite{Ref23}. Therefore, when a systems engineer utilizes set theory to define a complex architecture, they must construct a massive apparatus of external logical inference rules to reason about the sets they have defined \parencite{Ref23}. The theory itself contains axioms, but the manipulation of those axioms requires a separate logical framework \parencite{Ref16, Ref26}. This material ontology heavily dictates how dynamic relationships are modeled. In set theory, a function is not a primitive computational action; it is a static structural entity defined as a subset of the Cartesian product of two sets \parencite{Ref23}. Specifically, a function $f: A \to B$ is a relation (a set of ordered pairs) that must satisfy the external properties of totality (for every element $x \in A$, there exists an element $y \in B$) and single-valuedness (if the ordered pairs $(x, y_1)$ and $(x, y_2)$ are in the relation, then $y_1$ must equal $y_2$) \parencite{Ref23, Ref23}. When modeling systems engineering phenomena---such as the trajectory of a system state over time---an engineer using set theory must continuously generate manual proofs verifying that the subsets they generate adhere strictly to these relational properties \parencite{Ref27}. Furthermore, equality in set theory is strictly extensional. Two sets are deemed equal if and only if they contain the exact same elements, regardless of how those sets were constructed or defined. Two functions are equal if their relational graphs are identical sets of ordered pairs \parencite{Ref23}. While this static, extensional view of equality is highly effective for abstract, non-constructive mathematical proofs, it poses significant, often insurmountable challenges for algorithmic evaluation and automated reasoning. Checking the extensional equivalence of infinite or highly complex sets is generally uncomputable, rendering set-theoretic models exceptionally difficult to verify using software \parencite{Ref15, Ref16, Ref17}. As mathematical practitioners note, defining even low-complexity notions in set theory requires complex, often convoluted set-theoretic coding that is difficult to work with syntactically \parencite{Ref17}.\subsection{Dependent Type Theory: Structural Ontology and Internal Logic}
Dependent Type Theory, initiated by Per Martin-Löf and subsequently refined through frameworks such as the Calculus of Constructions, operates on fundamentally different philosophical and mathematical axioms \parencite{Ref18, Ref19, Ref20}. Designed specifically to address the computational shortcomings of classical foundations, DTT abandons the global, material universe of sets in favor of a rigid, structural hierarchy of types \parencite{Ref28}. In DTT, an element (or term) does not exist independently; it is intrinsically and permanently bound to its type \parencite{Ref16, Ref18}. The expression $a : A$ denotes that the term $a$ is a structural inhabitant of the type $A$ \parencite{Ref18}. Because types are structural, it is syntactically invalid to ask if a term of one type belongs to another disjoint type \parencite{Ref18}. This foundational constraint eliminates entire classes of paradoxes and logical errors natively \parencite{Ref25}. The most profound divergence between the two paradigms lies in the treatment of logic. DTT operates on the Curry-Howard isomorphism, a mathematical principle demonstrating a direct correspondence between computer programs and mathematical proofs. Under this paradigm, known as ``propositions-as-types,'' logic is not an external apparatus bolted onto the theory; it is internal and native to the type system itself \parencite{Ref16, Ref23}. A mathematical proposition is represented as a type, and a proof of that proposition is simply an executable program (a term) that inhabits that type \parencite{Ref16}. Because of this internal logic, traditional logical connectives are derived directly from type formers \parencite{Ref23}. Implication maps to standard function types, conjunction maps to product types, and, crucially, universal quantification maps to dependent product types \parencite{Ref23, Ref29}. Therefore, in DTT, type-checking a system specification is mathematically identical to verifying a formal proof of its correctness \parencite{Ref22, Ref30}. If a system model compiles, it is logically consistent by definition.

The defining and most powerful feature of DTT is the dependency mechanism itself: a type's specification may vary based on a specific value \parencite{Ref19, Ref29}. Let $U$ denote a universe of types. For a family of types $B$ parameterized by a term $x : A$, the dependent product type (often called a Pi-type, written as $\prod_{x : A} B(x)$) represents functions that take an argument $x$ and return a value of a type that specifically depends on the evaluation of $x$ \parencite{Ref29}. In the context of systems engineering, this allows the mathematical framework to enforce complex, dynamic invariants natively. For instance, the type of a system's output vector can be defined to depend strictly on the precise state value from which the system transitioned \parencite{Ref28}. Dependent sum types (Sigma-types) similarly allow for the bundling of a value with a proof that the value satisfies a specific condition, effectively acting as an indexed family of sets \parencite{Ref19, Ref29}. Unlike the static relational subsets of ZFC, functions in DTT are primitive computational constructs defined by lambda abstraction \parencite{Ref23}. They possess inherent, executable computational behavior \parencite{Ref19}. Furthermore, DTT introduces a highly nuanced treatment of equality \parencite{Ref25}. It formally distinguishes between propositional equality (which requires an explicit proof term to demonstrate that two entities are equal) and definitional (or judgemental) equality \parencite{Ref25}. If two terms reduce to the exact same normal form via algorithmic computation (for example, the function $f(x) = 2 + 2$ and the value $4$), they are definitionally equal \parencite{Ref19}. The compiler's type checker accepts them as identical without requiring any manual proof steps or relational mapping \parencite{Ref19, Ref25}.\subsection{The ``Types Are Not Sets'' Paradigm}
The literature surrounding formal programming semantics has long debated the conceptual mapping between these two realms. As early as 1973, James H. Morris Jr. published the seminal paper ``Types are not Sets,'' an assertion heavily supported by researchers like Leslie Lamport and Lawrence Paulson \parencite{Ref31, Ref32, Ref33, Ref34, Ref35}. The core argument of this literature is that types represent the \textit{syntax} and structural grammar of a logical system, whereas sets represent one possible \textit{semantic model} or interpretation of that syntax \parencite{Ref24}. Dana Scott and Neel Krishnaswami expanded on this, noting that viewing types purely as sets is mathematically reductive. While the models of well-behaved type theories often include sets, they also include complex entities that are not sets, such as realizability interpretations necessary for compiling executable programs \parencite{Ref24}. Homotopy Type Theory (HoTT) further expands this expressive capability, allowing types to be interpreted as higher-dimensional infinity-groupoids or simplicial sets, providing a mathematical language for continuous transformations that is uniquely difficult to represent in standard ZFC without invoking ``evil'' strict semantic interpretations \parencite{Ref17}. Ultimately, as a foundational system, type theory was designed explicitly for programming language theory and constructive proof extraction, whereas set theory was designed to study classical independence phenomena and infinitary combinatorics \parencite{Ref17}. Applying set theory directly to the engineering of computational cyber-physical systems is therefore an application of the wrong mathematical tool, leading to massive frictional overhead.

\begin{table}[h]
\centering
\renewcommand{\arraystretch}{1.3}
\begin{tabularx}{\textwidth}{@{}l X X@{}}
\toprule
\textbf{Analytical Dimension} & \textbf{Classical Set Theory (ZFC)} & \textbf{Dependent Type Theory (DTT)} \\
\midrule
\textbf{Ontological Basis} & Material (elements exist independent of sets) & Structural (terms are intrinsically typed) \\
\textbf{Logical Framework} & External (first-order predicate calculus) & Internal (propositions as types via Curry-Howard) \\
\textbf{Function Definition} & Derived relations over Cartesian products & Primitive computational entities (lambda calculus) \\
\textbf{Equality Paradigm} & Extensional (sets with identical elements) & Definitional (computational reduction) and Propositional \\
\textbf{Computational Nature} & Descriptive, non-executable, abstract & Prescriptive, statically analyzable, executable \\
\textbf{Expressive Target} & Infinitary combinatorics, abstract mathematics & Constructive algorithms, program verification \\
\bottomrule
\end{tabularx}
\caption{Comparative breakdown of Set Theory vs. Type Theory paradigms.}
\end{table}

\section{Wymore's Tricotyledon Theory of System Design (T3SD)}

Wayne Wymore's Tricotyledon Theory of System Design (T3SD) stands as one of the most ambitious and comprehensive attempts to rigorously define the systems engineering lifecycle. Developed to replace the ambiguous natural-language specifications prevalent in the aerospace and defense industries, T3SD spans from requirements analysis to system coupling, modeling, simulation, and physical implementation \parencite{Ref11, Ref12, Ref14}. It proposes a ``completely rigorous mathematical theory of system design'' \parencite{Ref12}. Within the T3SD framework, a system is modeled fundamentally as a port automaton utilizing a set-theoretic foundation \parencite{Ref11}. Wymore meticulously defined various cotyledons (perspectives) of design, including the functional, buildability, and implementability perspectives, driven by rigorous definitions of system requirements \parencite{Ref36}. Wymore defined the System Design Requirement (SDR) as a complex set containing Input/Output Requirements (IOR), System Test Requirements (STR), Cost Requirements (CR), Performance Requirements (PR), Tradeoff Requirements (TR), and Technology Requirements (TYR) \parencite{Ref11, Ref14, Ref36}.\subsection{The Discrete System Model}
The central architectural artifact of T3SD is the discrete system model, which Wymore formally defined via a specific set-theoretic quintuple $Z = (S_Z, I_Z, O_Z, N_Z, R_Z)$ \parencite{Ref27}. In Wymore's formalization:
\begin{itemize}
 \item $S_Z$ represents a non-empty set of system states.
 \item $I_Z$ represents the set of all possible inputs to the system.
 \item $O_Z$ represents the set of all possible outputs from the system.
 \item $N_Z: S_Z \times I_Z \to S_Z$ represents the state transition function, detailing how the system evolves given a current state and a specific input.
 \item $R_Z: S_Z \to O_Z$ represents the readout (or output) function, detailing the output generated from any given state.
\end{itemize}

To capture realistic engineering concepts such as system closure (systems with no external inputs or outputs) and autonomous state transitions (systems that evolve over time purely based on internal dynamics without external stimuli), Wymore had to utilize complex set-theoretic workarounds \parencite{Ref13, Ref27}. For example, he defined empty sets as valid input domains and utilized partial functions that mapped from obscure subsets of Cartesian products to represent behaviors where inputs were not present \parencite{Ref27}.\subsection{System Coupling and Closure}
Beyond individual system definitions, Wymore introduced a robust mathematical algebra for system architecture through the concept of ``coupling recipes'' \parencite{Ref14, Ref36, Ref37}. A system coupling recipe mathematically defines how the output ports of a specific set of subsystems connect to the input ports of other subsystems to form a larger, complex entity \parencite{Ref13, Ref38}. A paramount property of T3SD is the theorem of ``closure under coupling'' \parencite{Ref36}. Wymore provided exhaustive proofs demonstrating that any valid coupling of discrete subsystems results in a new, holistic entity that itself perfectly satisfies the mathematical definition of a singular discrete system (the quintuple defined above) \parencite{Ref13, Ref36}. While this is immensely powerful conceptually, Wymore's reliance on classical set theory meant that proving these coupling closures, along with state trajectory bounds and time invariances, required hundreds of pages of dense, manual mathematical proofs \parencite{Ref37}. Every time an engineer altered a coupling recipe, the relational mapping of the subsets had to be manually verified to ensure it still constituted a valid mathematical function.

\section{Translating T3SD into Dependent Type Theory: Context for Correctness}

The migration of Wymore's definitions from classical set theory to Dependent Type Theory demonstrates the massive methodological advantages of the structural, computational approach \parencite{Ref27}. The provided Lean 4 file \texttt{Wymore.lean} and the accompanying Proof Comparison Report offer a faithful encoding of T3SD Definition 2.4 (Discrete Systems) and provide empirical context for how DTT alters the paradigm of correctness and provability \parencite{Ref23, Ref23}.\subsection{The Core System Structure in DTT}
In the Lean 4 type-theoretic encoding, a discrete system is formulated not as a tuple of generic relational sets, but as a rigidly typed, parameterized \texttt{structure} \parencite{Ref27}:

\begin{verbatim}
structure DiscreteSystem (SZ : Type) (IZ : Type) (OZ : Type) where
 sz_nonempty : Nonempty SZ
 NZ : SZ → Option IZ → SZ
 RZ : SZ → Option OZ
\end{verbatim}

This precise DTT encoding introduces immediate conceptual and computational clarifications over the set-theoretic approach \parencite{Ref27}:
\begin{enumerate}
 \item \textbf{Universes and Types}: The state space $S_Z$, input space $I_Z$, and output space $O_Z$ are elevated from subsets of a universal set to distinct, parameterized computational \texttt{Type} universes \parencite{Ref27}.\item \textbf{The \texttt{Option} Monad for Partiality and Autonomy}: In set theory, modeling a system that can take autonomous steps or modeling a ``closed'' system requires defining transition functions over complex sub-domains or allowing mathematically dangerous partial functions. DTT utilizes the \texttt{Option} inductive type (which evaluates to either \texttt{some x} or \texttt{none}). The transition function \texttt{NZ : SZ $\to$ Option IZ $\to$ SZ} elegantly encapsulates both driven transitions (receiving an input \texttt{some i}) and autonomous transitions (receiving \texttt{none}) without breaking the totality requirement of the type system \parencite{Ref27}.\item \textbf{Proof Inhabitation for Correctness}: The field \texttt{sz\_nonempty : Nonempty SZ} is an application of the Curry-Howard isomorphism. It requires the systems engineer to actively provide a mathematical proof object demonstrating that the state space is capable of instantiation before the system can even be compiled. This rigidly enforces correct-by-construction modeling, guaranteeing that impossible systems cannot exist in the codebase \parencite{Ref27, Ref39, Ref40}.\end{enumerate}

\subsection{Eradicating Relational Overhead: Time Scales and Function Spaces}
In Wymore's textbook, time is established as a discrete scale $IJS^{++}$ mapping to the set of natural numbers. To represent system trajectories over time, Wymore relies heavily on the Function Space set, $\text{FNS}(A, B)$ \parencite{Ref23, Ref23}. To mathematically prove that a given trajectory is a valid function, an engineer must manually verify Wymore's Definition A1.155: that the graph of the relation meets the strict criteria for totality and single-valuedness \parencite{Ref27}. In DTT, this immense relational overhead is completely eradicated. Time is mapped directly to the native \texttt{Nat} (natural numbers) type \parencite{Ref23, Ref23}. The Lean implementation includes a theorem (\texttt{satisfiesFNS\_of\_function}) that computationally proves that \textit{any} natively defined function in Lean automatically satisfies Wymore's relational FNS criteria \parencite{Ref27}. Because DTT functions are primitives that must pass the compiler's strict termination and totality checkers to exist, the entire class of set-theoretic well-formedness proofs required by T3SD is bypassed \parencite{Ref23, Ref23}.\section{Proof Dynamics: The Fragility of Manual Deduction vs. Automated Verification}

The comparative analysis of textbook set theory versus Lean 4 type theory---as documented in the provided Proof Comparison Report \parencite{Ref27}---illuminates the transformative power of DTT in formalizing systems engineering mathematics. By analyzing the proofs of Wymore's core theorems, the divergence in epistemological friction becomes starkly evident.

\subsection{Theorem 2.25: Closure of Trajectories Under Translation and Concatenation}
Wymore postulates that complete input trajectories are closed under temporal translation (shifting time forward) and concatenation (joining two distinct trajectories at a specific time $r$) \parencite{Ref27}.\begin{itemize}
 \item \textbf{The Set-Theoretic Approach}: Wymore must utilize external theorems (A1.286 and A1.292) regarding set domains. To prove closure under translation for a function $f_1$ shifted by $t$, he must mathematically prove that the new domain $W' = \{s \mid s + t \in \text{TZ}\}$ is exactly equivalent to the original time scale $\text{TZ}$. For concatenation, he must compute the complex union of partitioned temporal intervals $\text{TZ}$.
 \item \textbf{The Type-Theoretic Approach}: In Lean, a complete trajectory is simply a total function of type \texttt{Time $\to$ A}. The \texttt{translate} and \texttt{concatenate} operations are defined directly as algorithmic functions. The type checker statically verifies that the return type of these algorithms is inherently \texttt{Time $\to$ A}. Closure is verified instantly by the compiler during type-checking without requiring any explicit proofs of domain shifting, set bounding, or set union \parencite{Ref27}.\end{itemize}

\subsection{Theorem 2.29: State Trajectory is a Function (Uniqueness)}
\begin{itemize}
 \item \textbf{The Set-Theoretic Approach}: Wymore defines the state trajectory mathematically as a subset of the Cartesian product $\text{TZ} \times \text{SZ}$. He then executes a manual, multi-page inductive proof to establish totality (a state exists for every time $t$) and single-valuedness (assuming the states $y_1 = y_2$ at time $t$, proving $y_1 = y_2$ at time $t+1$) \parencite{Ref27}.\item \textbf{The Type-Theoretic Approach}: Lean defines the trajectory simply via a recursive function (\texttt{generateStateTrajectory}). The compiler's recursion engine natively guarantees totality and well-formedness \parencite{Ref27}. Uniqueness is proven via the theorem \texttt{stateTrajectory\_unique}, utilizing structural induction on the time variable \texttt{t}. The burden of proof shifts fundamentally: instead of proving \textit{whether} the artifact is a valid mathematical object, the engineer only proves \textit{behavioral equivalence} to the generated baseline \parencite{Ref27}.\end{itemize}

\subsection{Theorem 2.32 and Erratum Discovery: The Vulnerability of Human Logic}
Perhaps the most crucial methodological insight derived from the comparative analysis is the discovery of a profound logical error in Wymore's original text \parencite{Ref27}.\begin{itemize}
 \item \textbf{The Set-Theoretic Failure}: To prove that the output trajectory is a valid function composed of the state trajectory and the readout function, Wymore explicitly cited Theorem A1.249 (image of preimage inclusion) and A1.250 (preimage of a complement) to justify his composition \parencite{Ref27}. However, neither theorem relates to function composition. This oversight illustrates the primary vulnerability of set-theoretic MBSE: human engineers are uniquely poor at managing massive, interdependent relational databases of axioms and lemmas without making semantic errors \parencite{Ref27}.\item \textbf{The Type-Theoretic Remedy}: In the Lean DTT encoding, output generation is defined by composing the readout function \texttt{Z.RZ} (type \texttt{SZ $\to$ Option OZ}) with the generated state trajectory (type \texttt{Time $\to$ SZ}). The Lean type checker automatically verifies the composite type as \texttt{Time $\to$ Option OZ}. Furthermore, the proof of pointwise identity is established via a single reflexivity command (\texttt{rfl}) because the equation is definitionally true by computation \parencite{Ref19, Ref27, Ref25}. DTT makes cross-referencing semantic errors logically impossible, as the compiler enforces absolute systemic consistency throughout the entire architecture \parencite{Ref39, Ref40}.\end{itemize}

\subsection{Theorem 2.46: Time Invariance of State Trajectory}
Wymore's proof of time invariance---that a trajectory shifted by $s$ and evaluated at $t$ is equal to the original trajectory evaluated at $s+t$---required a highly complex, multi-page double induction across two distinct variables \parencite{Ref27}. In stark contrast, the Lean 4 implementation achieves this proof with a single, highly compressed structural induction on $t$. Because DTT intrinsically understands the algebraic properties of its foundational types (utilizing \texttt{Nat.add\_comm} for the commutativity of addition), the proof is drastically simplified, completely mechanized, and devoid of the logical labyrinth required by set theory \parencite{Ref27}.\begin{table}[h]
\centering
\renewcommand{\arraystretch}{1.3}
\begin{tabularx}{\textwidth}{@{}l X X@{}}
\toprule
\textbf{Analytical Dimension} & \textbf{Set-Theoretic MBSE (T3SD)} & \textbf{Dependent Type-Theoretic MBSE (Lean/Coq)} \\
\midrule
\textbf{Verification Method} & Manual human deduction and cross-referencing & Automated compiler static type-checking \\
\textbf{Function Domains} & Explicit manual proofs of subset unions and bounds & Handled implicitly by domain typing and universes \\
\textbf{Composition Proofs} & Requires external relational lemmas (highly error-prone) & Definitionally true; proven trivially via reflexivity \\
\textbf{Error Vulnerability} & Extremely High (e.g., Theorem 2.32 cross-reference erratum) & Zero (compiler enforces absolute logical coherence) \\
\textbf{Inductive Reasoning} & Complex, multi-variable nested induction strategies & Simplified structural induction over inductive data types \\
\bottomrule
\end{tabularx}
\caption{Comparison of Set-Theoretic and Dependent Type-Theoretic Verification Dynamics.}
\end{table}

\section{Methodological and Conceptual Enhancements to Systems Engineering}

The translation of systems theory from classical set theory to dependent type theory offers far more than syntactic convenience; it provides sweeping methodological enhancements for the entire discipline of systems engineering. These enhancements align perfectly with the overarching mandates of the INCOSE FuSE initiative to ground the discipline in rigorous, computable science \parencite{Ref8, Ref9}.\subsection{1. Correct-by-Construction System Architectures}
Traditional software and systems engineering rely heavily on \textit{post hoc} testing to verify that a designed system meets its functional requirements \parencite{Ref41, Ref42}. MBSE methodologies that use informal graphical diagrams (like SysML) suffer from semantic ambiguity; the translation from an abstract diagram to a physical implementation often introduces fatal flaws that are only discovered during integration testing \parencite{Ref4, Ref43}. DTT fundamentally alters this lifecycle by enabling a ``correct-by-construction'' methodology \parencite{Ref6, Ref39, Ref40}. Because the dependent type system requires rigorous proofs of system properties to be passed as arguments before compilation is permitted, the design artifacts are mathematically guaranteed to satisfy their behavioral and structural constraints at the moment of specification \parencite{Ref22, Ref27}. A system that violates its safety constraints simply cannot be expressed or compiled in the language.

\subsection{2. Bridging the Gap Between Specification and Execution}
Set-theoretic models are generally descriptive and highly abstract; they cannot be directly executed or ``run'' by a machine \parencite{Ref15}. A T3SD state space model must be manually interpreted and translated into a distinct programming language (like C++ or Python) for simulation, breaking the digital thread. DTT, conversely, is prescriptive and executable \parencite{Ref19}. As demonstrated in the context provided, functions generating state trajectories are not merely theoretical constructs; they are directly executable code. Systems engineers can run interactive simulations of their formal models instantly within the same environment used for specification \parencite{Ref19}. This mathematically unifies the previously distinct activities of formal specification, modeling, and simulation into a singular, unbroken computational continuum \parencite{Ref22, Ref44}.\subsection{3. Managing Coupling Complexity via Polynomial Functors}
System coupling is the architectural core of complex systems engineering \parencite{Ref14, Ref36}. While Wymore utilized dense set-theoretic mapping to define his coupling recipes, modern category theory and type theory offer significantly superior mathematical tools. Recent advances in categorical systems theory utilize \textit{polynomial functors} and \textit{lenses} to describe interaction and dynamical systems algebraically \parencite{Ref45, Ref46, Ref47}. Polynomial functors, which integrate seamlessly with dependent type theories, allow engineers to rigorously model mode-dependent, open dynamical systems where the system's architecture, interfaces, and rewiring capabilities evolve dynamically over time \parencite{Ref46, Ref47}. This provides a robust, compositional mathematical foundation for the coupling of modern cyber-physical systems, overcoming the static, brittle rigidities of older port automaton models \parencite{Ref45, Ref48, Ref49}.\subsection{4. Extending Verification to Continuous and Hybrid Systems}
While Wymore's core formulation focuses heavily on discrete time and state spaces, modern cyber-physical systems---such as autonomous vehicles and power grids---operate in continuous environments governed by differential equations \parencite{Ref50, Ref51, Ref52}. The type-theoretic paradigm scales gracefully to continuous domains. Specifically, \textit{Temporal Type Theory} (developed extensively by David Spivak and Patrick Schultz) introduces a topos-theoretic approach to system behavior, providing a higher-order logic capable of internally defining derivatives and continuous real-valued functions without resorting to classical limits \parencite{Ref53, Ref54, Ref55, Ref56}. By embedding these theories into proof assistants, systems engineers can formally specify and verify complex hybrid dynamical systems, such as aircraft separation protocols within a national airspace---tasks that are virtually impossible to mechanize robustly using pure set theory \parencite{Ref55, Ref57}.\subsection{5. Semantic Retrofitting of Existing MBSE Frameworks}
The global systems engineering community is deeply invested in SysML and existing MBSE infrastructure \parencite{Ref5, Ref42, Ref58}. Recognizing this, the transition to formal methods does not require discarding these tools; rather, type theory offers a mechanism for semantic retrofitting. Researchers such as Henson Graves have successfully demonstrated how to map fragments of SysML structural block diagrams directly into type-theoretic semantics using Web Ontology Language (OWL) and formal logic constructs \parencite{Ref4, Ref59}. By defining an Abstract Block Diagram (ABD) as a mathematical theory closed under specific type constructions, the informal boxes and lines of SysML are translated into rigorous formalisms \parencite{Ref4, Ref60}. This allows engineers to continue using familiar graphical interfaces for conceptual design while leveraging the immense computational verification power of a type-theoretic backend to automatically check constraint violations, trace capability requirements, and ensure logical satisfiability across the architecture \parencite{Ref60, Ref61, Ref62}. Furthermore, as the industry explores integrating artificial intelligence into systems engineering, type theory provides the necessary structured ground truth. Frameworks utilizing Large Language Models (LLMs) to generate system architectures from natural language require formal modeling tools (like X Language or DTT backends) to ensure that the generated designs are logically consistent and executable, preventing the hallucination of impossible system states \parencite{Ref63, Ref64, Ref65}.\section{Proposed Research Angles for Leading Academic Journals}

To disseminate these findings, advance the theoretical frontier of the discipline, and directly influence the trajectory of the INCOSE FuSE initiative, the following comprehensive paper proposals are formulated. These proposals specifically target high-impact journals in the systems engineering domain, focusing on the profound methodological benefits of type theory.

\subsection{Angle 1: Target Journal: \textit{Systems Engineering} (Wiley)}
\noindent \textbf{Proposed Title:} From Relational Axioms to Computational Types: Re-founding MBSE on Dependent Type Theory\\
\textbf{Strategic Alignment:} This paper will directly address the INCOSE FuSE ``Foundations'' stream's mandate to rigorously evaluate and evolve the theoretical basis of systems engineering \parencite{Ref3, Ref8, Ref66}.\\
\textbf{Theoretical Framework and Argument:} The research will argue that the discipline's historical reliance on set-theoretic formulations (such as T3SD) fundamentally limits the efficacy, scalability, and verification potential of modern computational MBSE tools \parencite{Ref1, Ref8}. By structurally contrasting Wymore's set-theoretic framework with a Dependent Type Theory encoding, the paper will demonstrate how type theory mathematically unifies specification, simulation, and computation \parencite{Ref22}. The narrative will center heavily on the methodological superiority of the Curry-Howard isomorphism within engineering applications, showcasing how ``correct-by-construction'' paradigms and definitional equality drastically reduce verification burdens in complex system design \parencite{Ref6, Ref25, Ref39}. The paper will intentionally avoid framing the topic as purely a ``software implementation'' tutorial, instead positioning DTT philosophically and mathematically as the necessary underlying ontology for next-generation systems engineering \parencite{Ref1, Ref66}.\subsection{Angle 2: Target Journal: \textit{Systems Engineering} (Wiley)}
\noindent \textbf{Proposed Title:} Exposing Latent Vulnerabilities in Formal Systems Theory via Machine-Checked Proofs\\
\textbf{Strategic Alignment:} This paper leverages the errata discovered in foundational texts (specifically Wymore's Theorem 2.32) as a highly impactful case study demonstrating the fragility of manual mathematical models in safety-critical systems \parencite{Ref7, Ref27}.\\
\textbf{Theoretical Framework and Argument:} The research will explore how manual, set-theoretic proofs of system architectures are highly prone to human semantic error---specifically cross-referencing failures and logic cascade breakdowns \parencite{Ref27}. The paper will vividly contrast this fragility with the absolute logical coherence enforced by the compiler of a dependent type theory-based proof assistant \parencite{Ref21, Ref27}. By emphasizing the critical necessity of formal methods in aerospace, defense, and autonomous systems, this angle will demonstrate that DTT proof assistants are not merely academic novelties, but critical risk mitigation engines necessary for verifying system coupling, state bounds, and trajectory invariances in high-stakes environments \parencite{Ref7, Ref39, Ref59, Ref67}.\subsection{Angle 3: Target Journal: \textit{Systems} (MDPI)}
\noindent \textbf{Proposed Title:} A Type-Theoretic Ontology for Cyber-Physical Systems: Bridging SysML and Formal Generative Behavior\\
\textbf{Strategic Alignment:} This paper aligns perfectly with MDPI \textit{Systems}' scope regarding complex socio-technical systems, information systems, and the integration of emerging AI/LLM technologies with systems practice \parencite{Ref42, Ref64, Ref65}.\\
\textbf{Theoretical Framework and Argument:} Building upon the foundational work of retrofitting SysML with type-theoretic semantics \parencite{Ref4}, this paper will propose a unified modeling methodology that maps informal operational needs into executable DTT structures \parencite{Ref68}. It will highlight how type theory can capture rigid structural constraints while simultaneously embedding dynamic behavioral parameters natively. Furthermore, the paper will introduce how categorical extensions, specifically polynomial functors, can be used within a type-theoretic framework to manage dynamic interface complexity, generative design, and mode-dependent re-wiring in autonomous systems interacting with human agents \parencite{Ref45, Ref47, Ref49}. It will explicitly position DTT as the safest backend for LLM-assisted MBSE architectures \parencite{Ref65}.\subsection{Angle 4: Target Journal: \textit{Systems} (MDPI)}
\noindent \textbf{Proposed Title:} Executable Specifications in Hybrid Architectures: A Temporal Type-Theoretic Approach to Trajectory Validation\\
\textbf{Strategic Alignment:} Highlighting the practical application of formalized system trajectories in evolving, continuous environments, such as autonomous transport or dynamic risk assessment networks \parencite{Ref64}.\\
\textbf{Theoretical Framework and Argument:} Utilizing the modeling of state and output trajectories natively generated through recursive functions in DTT, this paper will demonstrate how time-invariant properties and non-anticipatory behaviors are mathematically guaranteed via structural induction rather than exhaustive simulation \parencite{Ref27}. The paper will abstract from discrete computational implementations to discuss continuous domains, utilizing Spivak's \textit{Temporal Type Theory} to show how guaranteed trajectories can be embedded into dynamic decision-support systems for cyber-physical architectures \parencite{Ref50, Ref55, Ref64}. The core narrative will focus on replacing late-stage post-design simulation validation with in-design continuous trajectory proof synthesis, drastically accelerating the verification lifecycle while ensuring mathematically flawless operations \parencite{Ref40, Ref56, Ref69}.\section{Conclusion}

The evolution of systems engineering from heuristic, document-centric practices to rigorous, model-based architectures demands a mathematical foundation commensurate with the exponential complexity of modern cyber-physical systems. While Wayne Wymore's Tricotyledon Theory of System Design provided a monumental theoretical leap by framing systems engineering within the rigorous boundaries of classical set theory, its fundamental reliance on relational structure, external material logic, and extensional equality renders it practically uncomputable for automated, software-driven verification. The inherent fragility of manual set-theoretic proofs---vividly evidenced by the logical errata exposed through machine-checking---demonstrates unequivocally that classical set theory is a descriptive language suited for abstract mathematics, not a prescriptive engineering tool suited for complex system architecture.

Dependent Type Theory provides the comprehensive epistemological and computational solution to the MBSE foundations crisis currently being addressed by the INCOSE FuSE initiative. By adopting a structural mathematical ontology where propositions are types, logical constraints are internal, and specifications act as executable algorithms, DTT completely eliminates the immense relational overhead required by classical methods. It replaces error-prone human deduction with infallible static compiler verification, enabling truly correct-by-construction system architectures where design flaws are caught at compilation rather than integration. 

Furthermore, the natural extensions of type theory into Categorical Systems Theory via polynomial functors, and into continuous domains via Temporal Type Theory, prove that this paradigm is not merely a software engineering construct, but a universal language for modeling complex, dynamic, mode-dependent interactions. As the systems engineering discipline seeks the future bedrock of its practice, the transition from set-theoretic models to type-theoretic paradigms represents not just a syntactic upgrade, but a fundamental, necessary leap in the human capacity to reliably specify, couple, verify, and execute the behavior of the world's most complex systems.


\section{References}
\printbibliography[heading=none]

\end{document}